\section{独立和、强弱大数律与中心极限}\label{sec:05-04}

\subsection{矩母函数、Chernoff 方法与集中不等式}\label{subsec:5-4-1}
\index{矩母函数}\index{Chernoff 方法}\index{集中不等式}

Chebyshev 不等式已经能回答一个基本问题。若 $X_1,\dots,X_n$ 独立同分布，
$\E X_i=\mu$、$\Var(X_i)=\sigma^2<\infty$，则
\[
 \Prob\left(\left|\frac1n\sum_{i=1}^nX_i-\mu\right|\geq\varepsilon\right)
 \leq \frac{\sigma^2}{n\varepsilon^2}.
\]
这个上界随 $n$ 衰减，却只按 $n^{-1}$ 衰减。对有界独立变量，实际尾概率通常按
$e^{-cn}$ 衰减。缺失的指数从哪里来？把和 $S$ 直接放进 Markov 不等式只会使用一阶矩；
把 $S$ 先送进单调函数 $x\mapsto e^{sx}$，则独立性会把和的指数变成可分解的乘积。
这一小步就是 Chernoff 方法的核心机制。

Chernoff 的指数变换、Hoeffding 的有界差估计与 Bernstein 的方差信息修正都遵循同一逻辑。
它们的共同点不是某个需背诵的公式，而是“指数化、因子化、
控制单项、优化参数”这组通用证明步骤。

在正式推导之前，先看三种截然不同的指数矩行为。

\begin{example}[Rademacher]\label{ex:5-4-1-1}
设 $X$ 以相同概率取 $1$ 与 $-1$。则 $\E X=0$，并且对每个 $t\in\R$，
\[
 \E e^{tX}=\frac{e^t+e^{-t}}2=\cosh t.
\]
令 $h(t)=\log\cosh t-t^2/2$。由于
$h'(t)=\tanh t-t$，而当 $t\geq0$ 时
$h''(t)=\operatorname{sech}^2t-1\leq0$，故 $h'(t)\leq h'(0)=0$，从而
$h(t)\leq h(0)=0$。函数 $h$ 为偶函数，所以
\[
 \cosh t\leq e^{t^2/2}\qquad(t\in\R).
\]
这正是下文定义的最基本的次高斯性质。
\end{example}

\begin{example}[高斯分布]\label{ex:5-4-1-2}
设 $G\sim N(0,\sigma^2)$，其中 $\sigma>0$。只在实轴上配方，便有
\begin{align*}
 \E e^{tG}
 &=\frac1{\sqrt{2\pi\sigma^2}}
   \int_{\R}\exp\left(tx-\frac{x^2}{2\sigma^2}\right)\dd x\\
 &=e^{\sigma^2t^2/2}\frac1{\sqrt{2\pi\sigma^2}}
   \int_{\R}\exp\left(-\frac{(x-\sigma^2t)^2}{2\sigma^2}\right)\dd x
 =e^{\sigma^2t^2/2}.
\end{align*}
第二个等号只是实变量平移 $u=x-\sigma^2t$；这里没有复变换元。高斯分布使
下文的次高斯上界取等，是以后比较常数的基准。
\end{example}

\begin{example}[Pareto 重尾变量]\label{ex:5-4-1-3}
令 $X$ 的密度为
\[
 p(x)=\alpha x^{-\alpha-1}\1_{[1,\infty)}(x),\qquad \alpha>0.
\]
首先
\[
 \int_1^\infty\alpha x^{-\alpha-1}\dd x
 =\bigl[-x^{-\alpha}\bigr]_{1}^{\infty}=1,
\]
所以这确实是概率密度。且对 $x\ge1$，
\[
 \Prob(X>x)=\int_x^\infty\alpha u^{-\alpha-1}\dd u=x^{-\alpha}.
\]
对任意 $t>0$，
\[
 \E e^{tX}
 \geq\alpha\sum_{n=1}^{\infty}
       \int_n^{n+1}e^{tx}x^{-\alpha-1}\dd x
 \geq\alpha\sum_{n=1}^{\infty}e^{tn}(n+1)^{-\alpha-1}=\infty,
\]
因为最后一级数的通项甚至不趋于零。因此正指数矩处处发散，Chernoff 方法在这里无从启动。
这不是计算技巧不足，而是重尾分布本身不具有有限的正指数矩；此时应改用截断或幂矩估计。
\end{example}

\begin{definition}[矩母函数与累积量生成函数]\label{def:5-4-1-1}
实值随机变量 $X$ 的矩母函数定义为
\[
 M_X(t)=\E e^{tX}\in(0,\infty],\qquad t\in\R.
\]
在 $M_X(t)<\infty$ 的参数集合上，记
$\Lambda_X(t)=\log M_X(t)$，称为\emph{累积量生成函数}。
在某一个正参数处有限已经足以给相应的单侧 Chernoff 界；要求零点邻域内处处有限是更强的条件，
二者不作混同。
\end{definition}

\begin{definition}[次高斯变量]\label{def:5-4-1-2}
\index{次高斯变量}
若实值随机变量 $X$ 满足 $\E X=0$，并且存在 $\sigma^2\geq0$ 使
\[
 \E e^{tX}\leq e^{\sigma^2t^2/2}\qquad(t\in\R),
\]
则称 $X$ 是参数 $\sigma^2$ 的\emph{次高斯}（sub-Gaussian）变量。这个参数是指数矩的控制常数，
并不按定义等于 $\Var(X)$。
\end{definition}

\begin{definition}[Chernoff 变换]\label{def:5-4-1-3}
设随机变量 $S$ 的对数矩母函数为 $\Lambda_S$。对 $a\in\R$，定义单侧 Chernoff 变换
\[
 I_S^+(a)=\sup_{\substack{s>0\\\Lambda_S(s)<\infty}}
 \{sa-\Lambda_S(s)\}.
\]
若上确界所取集合为空，则约定 $I_S^+(a)=-\infty$。
\end{definition}

\begin{proposition}[Chernoff 上界]
若至少存在一个 $s>0$ 使 $\E e^{sS}<\infty$，则
\[
 \Prob(S\geq a)
 \leq\inf_{\substack{s>0\\\Lambda_S(s)<\infty}}
 e^{-sa}\E e^{sS}
 =e^{-I_S^+(a)}.
\]
\end{proposition}

\begin{proof}
对每个可行的 $s>0$，指数函数的单调性和 Markov 不等式给出
\[
 \Prob(S\geq a)
 =\Prob(e^{sS}\geq e^{sa})
 \leq e^{-sa}\E e^{sS}.
\]
对 $s$ 取下确界，再用 $\Lambda_S(s)=\log\E e^{sS}$，便得到结论。
\end{proof}

真正需要估计的是单个有界中心变量的指数矩。凸性用两个端点的凸组合控制区间内的指数函数，而端点计算仍需一个
不会把最优常数藏起来的微积分步骤。

\begin{lemma}[Hoeffding 引理]\label{lemma:5-4-1-1}
设 $a\leq X\leq b$ a.s. 且 $\E X=0$。则对每个 $t\in\R$，
\[
 \E e^{tX}\leq\exp\left(\frac{t^2(b-a)^2}{8}\right).
\]
\end{lemma}

\begin{proof}
若 $a=b$，则 $X=0$ a.s.，结论成立。以下设 $a<b$。由 $\E X=0$ 可知
$a\leq0\leq b$。凸函数 $x\mapsto e^{tx}$ 在区间端点之间低于弦，因而对
$x\in[a,b]$，
\[
 e^{tx}\leq \frac{b-x}{b-a}e^{ta}+\frac{x-a}{b-a}e^{tb}.
\]
取期望并用 $\E X=0$，得到
\[
 \E e^{tX}\leq\frac{b}{b-a}e^{ta}-\frac{a}{b-a}e^{tb}.
\]
置
\[
 p=\frac{-a}{b-a}\in[0,1],\qquad s=t(b-a).
\]
右端变为
\[
 (1-p)e^{-ps}+pe^{(1-p)s}.
\]
令
\[
 H(s)=\log\bigl((1-p)e^{-ps}+pe^{(1-p)s}\bigr).
\]
把分母中的两项归一化为概率 $1-q_s,q_s$，其中
\[
 q_s=\frac{pe^{(1-p)s}}
 {(1-p)e^{-ps}+pe^{(1-p)s}},
\]
直接求导得到
\[
 H'(s)=q_s-p,\qquad H''(s)=q_s(1-q_s)\leq\frac14.
\]
又 $H(0)=H'(0)=0$。当 $s\geq0$ 时，Taylor 积分余项给
\[
 H(s)=\int_0^s(s-u)H''(u)\dd u\leq\frac{s^2}{8};
\]
当 $s<0$ 时，同一恒等式写成
$H(s)=\int_s^0(u-s)H''(u)\dd u$，仍得到相同上界。因此
$H(s)\leq s^2/8$ 对所有实数 $s$ 成立。代回 $s=t(b-a)$ 即得结论。
\end{proof}

独立性现在把每一项的估计相乘起来。

\begin{theorem}[Hoeffding 不等式]\label{theorem:5-4-1-2}
设 $X_1,\dots,X_n$ 独立，且
$a_i\leq X_i\leq b_i$ a.s.。则对每个 $u>0$，
若 $V:=\sum_{i=1}^n(b_i-a_i)^2>0$，则
\[
 \Prob\left(\sum_{i=1}^n(X_i-\E X_i)\geq u\right)
 \leq
 \exp\left(-\frac{2u^2}{V}\right).
\]
若 $V=0$，则左端为零。
\end{theorem}

\begin{proof}
记 $Y_i=X_i-\E X_i$ 以及
$V=\sum_i(b_i-a_i)^2$。若 $V=0$，则每个 $X_i$ a.s. 为常数，因而
$\sum_iY_i=0$ a.s.；结论成立。以下设 $V>0$。对任意 $s>0$，Chernoff 变换、
独立变量的期望因子化以及 Hoeffding 引理依次给出
\begin{align*}
 \Prob\left(\sum_iY_i\geq u\right)
 &\leq e^{-su}\E\exp\left(s\sum_iY_i\right)\\
 &=e^{-su}\prod_i\E e^{sY_i}\\
 &\leq\exp\left(-su+\frac{s^2}{8}
                  \sum_i(b_i-a_i)^2\right)
 =\exp\left(-su+\frac{s^2V}{8}\right).
\end{align*}
右端指数是关于 $s$ 的二次函数，在 $s=4u/V$ 处达到最小值
$-2u^2/V$。代入即得所述估计。
\end{proof}

\begin{corollary}[样本均值的双侧界]\label{corollary:5-4-1-hoeffding-two-sided}
若 $X_1,\dots,X_n$ 独立同分布，且 $a\leq X_i\leq b$ a.s.，其中 $a<b$，则
对每个 $\varepsilon>0$，
\[
 \Prob\left(
 \left|\frac1n\sum_{i=1}^nX_i-\E X_1\right|\geq\varepsilon
 \right)
 \leq2\exp\left(-\frac{2n\varepsilon^2}{(b-a)^2}\right).
\]
\end{corollary}

\begin{proof}
对上尾把定理中的 $u$ 取为 $n\varepsilon$。对变量 $-X_i$ 再应用一次定理得到下尾估计，
最后使用并集界。
\end{proof}

这个推论立即给出一个强结论。固定 $m\geq1$，右端在 $n$ 上可和，故 Borel--Cantelli
引理说明样本均值与期望相差至少 $1/m$ 的事件只发生有限次。对所有 $m$ 取可数交，得到
有界独立同分布变量的样本均值 a.s. 收敛到期望。这里不能只写“令 $\varepsilon\downarrow0$”：
把不可数多个满测集相交并不保测度一，可数阈值 $1/m$ 才完成这一步。

\begin{proposition}[次高斯变量之和]\label{proposition:5-4-1-3}
若 $X_1,\dots,X_n$ 独立、中心化，且 $X_i$ 是参数 $\sigma_i^2$ 的
次高斯变量，则 $\sum_iX_i$ 是参数 $\sum_i\sigma_i^2$ 的次高斯变量。
特别地，若 $v:=\sum_i\sigma_i^2>0$，则对每个 $u>0$，
\[
 \Prob\left(\sum_iX_i\geq u\right)
 \leq\exp\left(-\frac{u^2}{2v}\right).
\]
若 $v=0$，则 $\sum_iX_i=0$ a.s.
\end{proposition}

\begin{proof}
独立性与期望因子化给出
\[
 \E\exp\left(t\sum_iX_i\right)
 =\prod_i\E e^{tX_i}
 \leq\exp\left(\frac{t^2}{2}\sum_i\sigma_i^2\right).
\]
这就是第一个结论。若 $v=\sum_i\sigma_i^2>0$，Chernoff 上界的指数为
$-tu+t^2v/2$，取 $t=u/v$ 得尾界。若 $v=0$，则每个 $\sigma_i^2=0$，且
$\E e^{X_i}\leq1$。另一方面，$e^x\geq1+x$ 对所有实数 $x$ 成立，并且只在 $x=0$ 处取等，故
\[
 0\leq\E(e^{X_i}-1-X_i)\leq1-1-\E X_i=0.
\]
于是 $e^{X_i}-1-X_i=0$ a.s.，从而 $X_i=0$ a.s.。所以和为零 a.s.，所述上尾概率为零。
\end{proof}

Hoeffding 不等式只利用区间长度，没有利用实际方差。若方差远小于 $M^2$，下一条估计能把这份信息
放回指数中。

\begin{proposition}[Bernstein 型界（选读）]\label{proposition:5-4-1-4}
设 $X_1,\dots,X_n$ 为独立实值随机变量，$\E X_i=0$ 且
$|X_i|\leq M$ a.s.，其中 $M>0$。令
\[
 S=\sum_{i=1}^nX_i,\qquad v=\sum_{i=1}^n\E X_i^2.
\]
则对每个 $u>0$，
\[
 \Prob(S\geq u)
 \leq\exp\left(-\frac{u^2}{2(v+Mu/3)}\right).
\]
\end{proposition}

\begin{proof}
若 $0<s<3/M$，由 $\E X_i=0$、$|X_i|\leq M$ 及幂级数的绝对收敛，
\begin{align*}
 \E e^{sX_i}
 &=1+\sum_{k=2}^{\infty}\frac{s^k\E X_i^k}{k!}\\
 &\leq1+\sum_{k=2}^{\infty}
       \frac{s^kM^{k-2}\E X_i^2}{k!}.
\end{align*}
对 $k\geq2$ 有 $k!\geq2\,3^{k-2}$，所以
\[
 \E e^{sX_i}
 \leq1+\frac{s^2\E X_i^2}{2}
        \sum_{j=0}^{\infty}\left(\frac{sM}{3}\right)^j
 =1+\frac{s^2\E X_i^2}{2(1-sM/3)}.
\]
由 $\log(1+x)\leq x$，
\[
 \log\E e^{sX_i}
 \leq\frac{s^2\E X_i^2}{2(1-sM/3)}.
\]
若 $v=0$，则 $\E X_i^2=0$ 对每个 $i$ 成立，因而 $X_i=0$ a.s.，所述尾概率为零。
以下设 $v>0$。独立性与 Chernoff
变换给
\[
 \Prob(S\geq u)
 \leq\exp\left(-su+\frac{s^2v}{2(1-sM/3)}\right).
\]
取
$s=u/(v+Mu/3)$。此时 $sM/3=Mu/(3v+Mu)<1$。记 $D=v+Mu/3$，则
$s=u/D$ 且 $1-sM/3=v/D$，因而
\begin{align*}
 -su+\frac{s^2v}{2(1-sM/3)}
 &=-\frac{u^2}{D}
   +\frac{(u^2/D^2)v}{2(v/D)}\\
 &=-\frac{u^2}{D}+\frac{u^2}{2D}
 =-\frac{u^2}{2(v+Mu/3)}.
\end{align*}
\end{proof}

三种尾界的尺度现在可以直接比较。对 $[a,b]$ 中的 i.i.d. 变量，Chebyshev 给
$O((n\varepsilon^2)^{-1})$，Hoeffding 给
$2e^{-2n\varepsilon^2/(b-a)^2}$；Bernstein 在小偏差区近似
$e^{-n\varepsilon^2/(2\Var X_1)}$，在大偏差区转成 $e^{-cn\varepsilon/M}$。
例如若要以至少 $1-\delta$ 的置信度保证样本均值误差小于 $\varepsilon$，Hoeffding 给出充分条件
\[
 n\geq\frac{(b-a)^2}{2\varepsilon^2}\log\frac2\delta.
\]
这就是经验频率和有界独立抽样平均的有限样本保证。

\begin{center}
\begin{tikzpicture}[node distance=8mm and 9mm]
 \node[book strong node] (tail) {尾事件\\$\{S\geq u\}$};
 \node[book node,right=of tail] (exp) {指数化\\$e^{sS}\geq e^{su}$};
 \node[book node,right=of exp] (prod) {独立分解\\$\E e^{sS}=\prod_i\E e^{sX_i}$};
 \node[book warm node,below=of prod] (quad) {单项控制\\$\log\E e^{sX_i}\leq c_i(s)$};
 \node[book strong node,left=of quad] (opt) {参数优化\\$\inf_{s>0}$};
 \draw[book accent arrow] (tail)--(exp);
 \draw[book accent arrow] (exp)--(prod);
 \draw[book accent arrow] (prod)--(quad);
 \draw[book accent arrow] (quad)--(opt);
\end{tikzpicture}
\end{center}

图中的四步分别使用单调性、独立性、单变量估计与实变量优化。任何一步的假设失效，都不能靠
“集中现象应当成立”的直觉补齐。Pareto 例使第一项指数矩已经发散；相关变量使乘积公式失效；
次高斯参数也不是通常方差的另一个名字。

若失去独立性，上述因子化必须由针对具体依赖结构的估计取代。本小节的结论只针对独立变量。

\begin{remark}[本章只使用实参数矩母函数]\label{rem:5-4-1-complex-mgf}
本章不定义复参数矩母函数，也不调用其全纯性理论。Chernoff 估计只需要实数 $t>0$；
对 Rademacher 变量，实参数矩母函数是
$M(t)=(e^t+e^{-t})/2=\cosh t$，而 Pareto 变量对每个 $t>0$ 都有
$M(t)=\infty$。与之相反，第~\ref{subsec:5-4-3}~节所讨论的特征函数
$\E e^{itX}$ 因 $|e^{itX}|=1$ 而对每个实数 $t$ 都存在。
\end{remark}

\begin{exercise}[Bernoulli 平均的集中估计]
设 $B_1,\dots,B_n$ 为独立 Bernoulli$(p)$ 变量。
\begin{enumerate}
 \item 由 Hoeffding 不等式证明
 $\Prob(|n^{-1}\sum_iB_i-p|\geq\varepsilon)\leq2e^{-2n\varepsilon^2}$；
 \item 比较同一问题的 Chebyshev 上界，并求出两种方法为达到误差概率 $\delta$ 所给的样本量；
 \item 验证若去掉独立性并令 $B_i=B_1$，Hoeffding 的 $e^{-cn}$ 结论可能失败。
\end{enumerate}
\end{exercise}

\subsection{弱大数律、强大数律与截断方法}\label{subsec:5-4-2}
\index{弱大数律}\index{强大数律}\index{Kolmogorov 最大不等式}

集中不等式同时控制有限样本和渐近行为，但它要求有界性或指数矩。若只有二阶矩，
Chebyshev 不等式仍能给依概率收敛；要把结论升级为 a.s. 收敛，逐个时刻的概率趋零还不够，
必须控制一整段路径中的最大振荡。若进一步只保留一阶绝对矩，连方差都可能不存在，
于是还要先把罕见的大跳截掉。本小节依次完成这三次升级。

Bernoulli 的频率定理给出了大数律的早期形式；Borel 把结论推进到路径上的强收敛，Kolmogorov
又用最大不等式与截断给出自然的 $L^1$ 边界。下面的证明顺序保留这条技术升级：先缩放方差，
再控制二进分块，最后截去稀有大值。

\begin{definition}[样本和与样本平均]\label{def:5-4-2-1}
对随机变量序列 $(X_n)$，记
\[
 S_n=\sum_{k=1}^nX_k,
 \qquad
 \overline X_n=\frac{S_n}{n}.
\]
缩写 i.i.d. 表示这些变量相互独立且同分布。
\end{definition}

先从只需一行方差计算的结论开始。

\begin{theorem}[弱大数律（有限方差）]\label{theorem:5-4-2-1}
若 $(X_n)$ i.i.d.，$\E X_1=\mu$ 且 $\Var(X_1)=\sigma^2<\infty$，则
\[
 \frac{S_n}{n}\xrightarrow{\Prob}\mu.
\]
\end{theorem}

\begin{proof}
独立性和中心化给出
\[
 \Var\left(\frac{S_n}{n}\right)
 =\frac1{n^2}\sum_{k=1}^n\Var(X_k)
 =\frac{\sigma^2}{n}.
\]
因此对每个 $\varepsilon>0$，Chebyshev 不等式给
\[
 \Prob\left(\left|\frac{S_n}{n}-\mu\right|\geq\varepsilon\right)
 \leq\frac{\sigma^2}{n\varepsilon^2}\longrightarrow0.
\]
\end{proof}

这个证明不能直接给强律，因为 $\sum_n n^{-1}$ 发散，Borel--Cantelli 无法读取这组上界。
把时刻稀疏到 $2^k$ 后，上界变成可和；随之产生的新问题是如何控制相邻两个二进时刻之间的所有中间时刻。

\begin{example}[Bernoulli 频率]\label{ex:5-4-2-1}
设 $X_i=\1_{A_i}$ i.i.d.，且 $\Prob(A_i)=p$。由推论
\ref{corollary:5-4-1-hoeffding-two-sided}，
\[
 \Prob\left(\left|\frac{S_n}{n}-p\right|>\varepsilon\right)
 \leq2e^{-2n\varepsilon^2}.
\]
对每个 $\varepsilon=1/m$，右端在 $n$ 上可和。Borel--Cantelli I 和可数交于是给出
$S_n/n\to p$ a.s.。在这个有界模型中，集中估计已经给出了频率的强收敛。
\end{example}

可积性是一般强律的边界条件。标准 Cauchy 分布给出一个可计算的反例；其卷积稳定性将在
\cref{ex:5-4-3-cauchy-stability}中用纯实变量方法证明。

\begin{example}[Cauchy 反例]\label{ex:5-4-2-2}
标准 Cauchy 分布的密度为
\[
 f_1(x)=\frac1{\pi(1+x^2)}.
\]
归一化由
\[
 \int_{\R}f_1(x)\,\dd x
 =\frac1\pi\bigl[\arctan x\bigr]_{-\infty}^{\infty}=1
\]
得到。它不满足 $\E|X|<\infty$，因为
\[
 \int_1^R\frac{x}{\pi(1+x^2)}\,\dd x
 =\frac1{2\pi}\log\frac{1+R^2}{2}\longrightarrow\infty;
\]
由密度的对称性，正部和负部的期望
也都为无穷，所以 $\E X$ 不存在。例~\ref{ex:5-4-3-cauchy-stability} 将直接证明：若
$X_1,\dots,X_n$ 独立且均为标准 Cauchy，则
\[
 \frac1n\sum_{k=1}^nX_k\overset{\mathrm d}=X_1
\]
对每个 $n$ 成立。于是这些平均不可能依概率收敛到某个常数：若它们依概率收敛到 $c$，
便也依分布收敛到 $c$，但其分布恒为非退化 Cauchy 分布。这个例子说明删去可积性后强律可能失败；
它没有声称每个不可积分布都以同一种方式失败。
\end{example}

\begin{example}[有限方差的二进尺度]\label{ex:5-4-2-3}
令 $T_n=S_n-n\mu$。Chebyshev 在 $n=2^k$ 时给
\[
 \Prob(|T_{2^k}|>\varepsilon2^k)
 \leq\frac{\sigma^2}{\varepsilon^2 2^k},
\]
其和有限。故 $T_{2^k}/2^k\to0$ a.s.。然而这只控制二进时刻；
在区间 $[2^k,2^{k+1}]$ 内仍可能出现大幅偏离。下面的最大不等式正是填补这些区间的工具。
\end{example}

\begin{theorem}[Kolmogorov 最大不等式]\label{theorem:5-4-2-2}
设 $X_1,\dots,X_n$ 独立、实值、中心化且属于 $L^2$，并令
$S_k=\sum_{j=1}^kX_j$。则对每个 $\lambda>0$，
\[
 \Prob\left(\max_{1\leq k\leq n}|S_k|\geq\lambda\right)
 \leq\frac{\Var(S_n)}{\lambda^2}.
\]
\end{theorem}

\begin{proof}
对 $1\leq k\leq n$，令
\[
 A_k=\{|S_1|<\lambda,\dots,|S_{k-1}|<\lambda,
            |S_k|\geq\lambda\}.
\]
这些首次越界事件两两不交，其并为
$\{\max_{j\leq n}|S_j|\geq\lambda\}$。事件 $A_k$ 与 $S_k$ 都对
$\sigma(X_1,\dots,X_k)$ 可测，而增量 $S_n-S_k$ 与该 $\sigma$-代数独立并且均值为零。
由 Cauchy--Schwarz，不同因子的乘积可积；再用独立变量的期望因子化，得到
\[
 \E\bigl[S_k(S_n-S_k)\1_{A_k}\bigr]
 =\E[S_k\1_{A_k}]\,\E[S_n-S_k]=0.
\]
展开平方并使用最后一项非负，得到
\begin{align*}
 \E[S_n^2\1_{A_k}]
 &=\E[S_k^2\1_{A_k}]
   +2\E[S_k(S_n-S_k)\1_{A_k}]
   +\E[(S_n-S_k)^2\1_{A_k}]\\
 &\geq\E[S_k^2\1_{A_k}]
 \geq\lambda^2\Prob(A_k).
\end{align*}
对 $k$ 求和，并利用 $A_k$ 两两不交，
\[
 \lambda^2\Prob\left(\max_{j\leq n}|S_j|\geq\lambda\right)
 \leq\sum_{k=1}^n\E[S_n^2\1_{A_k}]
 \leq\E S_n^2=\Var(S_n).
\]
除以 $\lambda^2$ 即得结论。
\end{proof}

\begin{center}
\begin{tikzpicture}[x=1.05cm,y=0.8cm]
 \draw[->,book line] (0,0)--(9.5,0) node[right] {$k$};
 \draw[BookWarning,dashed] (0,2.6)--(9.1,2.6) node[right] {$+\lambda$};
 \draw[BookWarning,dashed] (0,-2.6)--(9.1,-2.6) node[right] {$-\lambda$};
 \draw[book accent arrow] (0,0)--(1,0.6)--(2,0.2)--(3,1.1)--(4,1.7)--(5,2.9)--(6,2.2)--(7,3.1)--(8,2.5)--(9,3.3);
 \fill[BookWarning] (5,2.9) circle (2pt);
 \draw[dashed,BookWarning] (5,0)--(5,2.9);
 \node[book warning node,above=3pt] at (5,2.9) {$A_5$：首次越界};
 \node[book node] at (7.2,0.9) {越界后的增量\\与过去独立且中心化};
\end{tikzpicture}
\end{center}

图只解释事件分割；交叉项为何消失仍由上面的可测性、独立性和条件期望计算保证。

\begin{theorem}[强大数律（有限方差）]\label{theorem:5-4-2-3}
若 $(X_n)$ i.i.d. 且 $X_1\in L^2$，则令 $\mu=\E X_1$，有
\[
 \frac{S_n}{n}\longrightarrow\mu\qquad\text{a.s.}
\]
\end{theorem}

\begin{proof}
置 $W_i=X_i-\mu$、$T_n=\sum_{i=1}^nW_i$，并记
$\sigma^2=\E W_1^2$。固定 $\varepsilon>0$。上一例已经给出
\[
 \sum_{k=0}^{\infty}
 \Prob(|T_{2^k}|>\varepsilon2^k)<\infty.
\]
另一方面，把最大不等式应用于独立块
$W_{2^k+1},\dots,W_{2^{k+1}}$，得到
\begin{align*}
 &\Prob\left(
 \max_{1\leq j\leq2^k}
 |T_{2^k+j}-T_{2^k}|>\varepsilon2^k
 \right)\\
 &\hspace{35mm}\leq
 \frac{2^k\sigma^2}{\varepsilon^2 2^{2k}}
 =\frac{\sigma^2}{\varepsilon^2 2^k}.
\end{align*}
这组概率同样可和。先令 $\varepsilon$ 遍历 $1/m$，再用 Borel--Cantelli I 并取可数交，
得到一个概率一事件；在该事件上，
\[
 \frac{|T_{2^k}|}{2^k}\longrightarrow0,
 \qquad
 \frac1{2^k}\max_{1\leq j\leq2^k}
 |T_{2^k+j}-T_{2^k}|\longrightarrow0.
\]
若 $2^k\leq n<2^{k+1}$，则
\[
 \frac{|T_n|}{n}
 \leq\frac{|T_{2^k}|}{2^k}
 +\frac1{2^k}\max_{1\leq j\leq2^k}
 |T_{2^k+j}-T_{2^k}|.
\]
右端趋于零，因此 $T_n/n\to0$，也就是 $S_n/n\to\mu$ a.s.
\end{proof}

有限方差并不是强律的最终边界。为了把条件降到 $L^1$，需要把一个随机级数的收敛转成
加权平均的收敛。下面的实分析引理承担这个转换。

\begin{lemma}[Kronecker 引理]\label{lemma:5-4-2-3}
设 $0<a_1\leq a_2\leq\cdots$ 且 $a_n\to\infty$。若实数或复数级数
$\sum_{n=1}^{\infty}x_n/a_n$ 收敛，则
\[
 \frac1{a_n}\sum_{k=1}^nx_k\longrightarrow0.
\]
\end{lemma}

\begin{proof}
令 $b_n=\sum_{k=1}^nx_k/a_k$、$b_0=0$，并设 $b_n\to b$。Abel 分部求和给
\[
 \sum_{k=1}^nx_k
 =\sum_{k=1}^na_k(b_k-b_{k-1})
 =a_nb_n-\sum_{k=1}^{n-1}(a_{k+1}-a_k)b_k.
\]
把 $b_k=b+(b_k-b)$ 代入并除以 $a_n$。常数 $b$ 的贡献为
$a_1b/a_n\to0$。再给定 $\varepsilon>0$，取 $K$ 使
$|b_k-b|<\varepsilon$ 对 $k\geq K$ 成立。有限前段的贡献除以 $a_n$ 后趋于零，而尾段满足
\[
 \frac1{a_n}\sum_{k=K}^{n-1}(a_{k+1}-a_k)|b_k-b|
 \leq\varepsilon\frac{a_n-a_K}{a_n}\leq\varepsilon.
\]
同时 $|b_n-b|\to0$。先令 $n\to\infty$，再令 $\varepsilon\downarrow0$，得到结论。
\end{proof}

\begin{theorem}[Kolmogorov 收敛判据]\label{theorem:5-4-2-kolmogorov-series}
设 $(Z_n)$ 独立、实值、中心化且 $Z_n\in L^2$。若
\[
 \sum_{n=1}^{\infty}\Var(Z_n)<\infty,
\]
则级数 $\sum_nZ_n$ 在 $L^2$ 中并且 a.s. 收敛。
\end{theorem}

\begin{proof}
令 $R_N=\sum_{n=1}^NZ_n$。独立性和中心化消去不同指标间的交叉项，所以对 $M>N$，
\[
 \E|R_M-R_N|^2=\sum_{n=N+1}^M\Var(Z_n)\longrightarrow0.
\]
$L^2$ 的完备性给出 $R_N$ 的 $L^2$ 极限。

为证明路径收敛，选取严格递增整数 $N_j$，使
\[
 \sum_{n>N_j}\Var(Z_n)\leq2^{-3j}.
\]
对每个 $j$，把 Kolmogorov 最大不等式用于有限块
$Z_{N_j+1},\dots,Z_{N_{j+1}}$，得到
\begin{align*}
 &\Prob\left(
 \max_{N_j<m\leq N_{j+1}}
 \left|R_m-R_{N_j}\right|>2^{-j}
 \right)\\
 &\hspace{30mm}\leq
 2^{2j}\sum_{n=N_j+1}^{N_{j+1}}\Var(Z_n)
 \leq2^{-j}.
\end{align*}
概率和有限，故 Borel--Cantelli I 表明：对 a.e. $\omega$，存在 $J(\omega)$，使对
$j\geq J(\omega)$，该块中的每个部分和都与块起点相差至多 $2^{-j}$。
若 $m>n\geq N_J$，设 $n,m$ 所在块的指标分别为 $j_0\leq j_1$。具体地，
\[
 R_m-R_n
 =-(R_n-R_{N_{j_0}})
  +\sum_{j=j_0}^{j_1-1}(R_{N_{j+1}}-R_{N_j})
  +(R_m-R_{N_{j_1}}),
\]
其中 $j_0=j_1$ 时中间的和为空和。两个端块增量与每个完整块增量都受相应的
$2^{-j}$ 控制，因而
\[
 |R_m(\omega)-R_n(\omega)|
 \leq2^{-j_0}+\sum_{j=j_0}^{j_1-1}2^{-j}+2^{-j_1}
 \leq4\,2^{-j_0}\leq4\,2^{-J}.
\]
让 $J\to\infty$，右端趋于零，所以 $(R_N(\omega))$ 是 Cauchy 数列并收敛。
记这个逐点极限为 $R_\infty$。前半部分已经给出某个 $R\in L^2$，使
$R_N\to R$ 于 $L^2$，从而依概率收敛；刚得到的 a.s. 收敛也蕴含
$R_N\xrightarrow{\Prob}R_\infty$。对任意 $\varepsilon>0$，
\[
 \Prob(|R-R_\infty|>\varepsilon)
 \leq
 \Prob(|R-R_N|>\varepsilon/2)
 +\Prob(|R_N-R_\infty|>\varepsilon/2)\longrightarrow0.
\]
因此 $R=R_\infty$ a.s.，这里的 $L^2$ 极限与路径极限确为同一个随机变量。
\end{proof}

作为抵消机制的一个例子，令 $\varepsilon_n$ 为独立 Rademacher 变量并取
$Z_n=\varepsilon_n/n$。方差级数 $\sum n^{-2}$ 收敛，故
$\sum\varepsilon_n/n$ a.s. 收敛；但
$\sum|\varepsilon_n|/n=\sum1/n$ 在每个样本点上发散。判据利用的是独立中心项的抵消，
不是绝对收敛的改写。

\begin{definition}[逐项截断]\label{def:5-4-2-2}
对 i.i.d. 序列 $(X_n)$，定义
\[
 Y_n=X_n\1_{\{|X_n|\leq n\}}.
\]
\end{definition}

\begin{lemma}[可积变量的逐项截断最终不变]\label{lemma:5-4-2-eventual-truncation}
若 $\E|X_1|<\infty$，则定义~\ref{def:5-4-2-2} 中的截断满足 $X_n=Y_n$ 最终 a.s.
\end{lemma}

\begin{proof}
若 $\E|X_1|<\infty$，则由 Tonelli 定理
\begin{align*}
 \sum_{n=1}^{\infty}\Prob(X_n\ne Y_n)
 &=\sum_{n=1}^{\infty}\Prob(|X_1|>n)
 =\E\sum_{n=1}^{\infty}\1_{\{|X_1|>n\}}\\
 &=\E\,\#\{n\in\N:n<|X_1|\}
 \leq\E|X_1|<\infty.
\end{align*}
因此 Borel--Cantelli I 保证 $X_n=Y_n$ 最终 a.s.。
\end{proof}

\begin{theorem}[Kolmogorov 强大数律]\label{theorem:5-4-2-4}
若 $(X_n)$ i.i.d. 且 $\E|X_1|<\infty$，则
\[
 \frac{S_n}{n}\longrightarrow\E X_1\qquad\text{a.s.}
\]
\end{theorem}

\begin{proof}
使用定义~\ref{def:5-4-2-2} 中的 $Y_n$。这些变量仍然独立。先验证随机级数判据的方差条件：
\begin{align*}
 \sum_{n=1}^{\infty}\frac{\Var(Y_n)}{n^2}
 &\leq\sum_{n=1}^{\infty}\frac{\E Y_n^2}{n^2}\\
 &=\E\left[|X_1|^2
       \sum_{\substack{n\geq1\\n\geq|X_1|}}\frac1{n^2}\right].
\end{align*}
存在绝对常数 $C$，使对每个 $r\geq0$，
\[
 r^2\sum_{\substack{n\geq1\\n\geq r}}\frac1{n^2}\leq Cr.
\]
事实上，$0\leq r<1$ 时用 $r^2\leq r$ 与 $\sum n^{-2}<\infty$；$r\geq1$ 时令
$N=\lceil r\rceil$，并用
$\sum_{n=N}^{\infty}n^{-2}\leq N^{-2}+\int_N^\infty x^{-2}\dd x\leq2/r$。
所以
\[
 \sum_{n=1}^{\infty}\frac{\Var(Y_n)}{n^2}
 \leq C\E|X_1|<\infty.
\]

令 $Z_n=(Y_n-\E Y_n)/n$。Kolmogorov 收敛判据说明
$\sum_nZ_n$ a.s. 收敛。对每条使该级数收敛的路径，把 Kronecker 引理用于
$a_n=n$、$x_n=Y_n-\E Y_n$，得到
\[
 \frac1n\sum_{k=1}^n(Y_k-\E Y_k)\longrightarrow0
 \qquad\text{a.s.}
\]
另一方面，
\[
 \E Y_n=\E\bigl[X_1\1_{\{|X_1|\leq n\}}\bigr]\longrightarrow\E X_1
\]
由支配收敛定理成立，故 \Cesaro{} 平均给
$n^{-1}\sum_{k\leq n}\E Y_k\to\E X_1$。

最后，\cref{lemma:5-4-2-eventual-truncation} 给出 $X_n=Y_n$ 最终 a.s.。在该概率一事件上，
$\sum_{k=1}^n(X_k-Y_k)$ 从某个时刻起是一个有限随机常数，除以 $n$ 后趋于零。
把三部分相加即得
\[
 \frac1n\sum_{k=1}^nX_k
 =\frac1n\sum_{k=1}^n(Y_k-\E Y_k)
  +\frac1n\sum_{k=1}^n\E Y_k
  +\frac1n\sum_{k=1}^n(X_k-Y_k)
 \longrightarrow\E X_1
\]
a.s.
\end{proof}

弱律与强律之间不是一个可以省略证明的箭头。弱律只控制每个固定时刻的坏事件；强律要求一条路径上
从某时刻起所有时刻都好。最大不等式把逐点控制升级为整块控制，截断再把罕见大跳交给
Borel--Cantelli。三种工具各自承担一段不同的逻辑。

因此经验均值、经验频率与独立抽样平均都由定理~\ref{theorem:5-4-2-4} 得到 a.s.
一致性。相关样本需要另行证明与其依赖结构相适应的极限定理；独立强律可以作为比较基准，却不能
直接用于相关样本。

\begin{exercise}[强律证明的关键步骤]
\begin{enumerate}
 \item 设 $(X_n)$ 独立、中心且 $\sum_n\Var(X_n)/n^2<\infty$。证明
 $n^{-1}\sum_{k\leq n}X_k\to0$ a.s.；
 \item 在有限方差强律的证明中，把二进时刻换成 $r^k$（固定 $r>1$），写出块最大估计；
 \item 构造一列 $X_n\xrightarrow{\Prob}0$ 但不 a.s. 收敛的变量，说明弱律形式本身不能推出强律。
\end{enumerate}
\end{exercise}

\subsection{特征函数、\Levy{} 连续性与中心极限定理}\label{subsec:5-4-3}
\index{特征函数}\index{Levy@\Levy{} 连续性定理}\index{中心极限定理}

大数律把 $S_n/n$ 的随机性压到一个常数。若要观察剩余波动，正确尺度通常是 $\sqrt n$。
此时直接计算归一化和的密度并不现实：每增加一项就要再做一次卷积，变量没有密度时这条路线甚至
无法起步。我们需要一种对每个概率分布都存在、又能把独立和变成乘积的观测。

矩母函数有乘法性质，却可能像 Pareto 例那样根本不存在。把实指数换成模长恒为一的振荡函数，
可积性问题消失了；代价是最后必须证明这种振荡观测足以唯一识别分布，并且其极限仍来自概率测度。
这两个缺口分别由唯一性定理与 \Levy{} 连续性定理填补。

de Moivre--Laplace 的二项近似早于现代测度论；\Levy{} 的贡献之一，是把这种计算组织成
普遍适用的特征函数连续性原理。现代证明必须同时交代两件事：紧性产生极限，唯一性识别极限。
把其中任何一件压缩掉，都会在质量逃逸的例子前失效。

\begin{definition}[特征函数]\label{def:5-4-3-1}
若 $X$ 是 $\R^d$ 值随机向量，其特征函数为
\[
 \varphi_X(t)=\E e^{i\langle t,X\rangle},\qquad t\in\R^d.
\]
若 $\mu$ 是 $\R^d$ 上的概率测度，也记
$\varphi_\mu(t)=\int e^{i\langle t,x\rangle}\mu(\dd x)$。
由于被积函数模长为一，特征函数对每个概率分布都有定义，并满足
$|\varphi_X(t)|\leq1$、$\varphi_X(0)=1$。
\end{definition}

\begin{definition}[正定函数]\label{def:5-4-3-2}
函数 $\varphi:\R^d\to\C$ 称为正定的，若对任意
$m\geq1$、$t_1,\dots,t_m\in\R^d$ 与 $c_1,\dots,c_m\in\C$，
\[
 \sum_{j,k=1}^mc_j\overline{c_k}\,\varphi(t_j-t_k)\geq0.
\]
这里左端按要求是非负实数。稍后证明每个特征函数都正定；正定连续函数何时反过来是特征函数，
属于 Bochner 定理的内容，这里不讨论该逆命题。
\end{definition}

\begin{proposition}[特征函数的基本性质]\label{proposition:5-4-3-1}
设 $X$ 为 $\R^d$ 值随机向量。
\begin{enumerate}
 \item $\varphi_X$ 在 $\R^d$ 上一致连续并且正定；
 \item 若 $X,Y$ 独立且同为 $\R^d$ 值，则
 $\varphi_{X+Y}=\varphi_X\varphi_Y$；
 \item 若 $A$ 是实 $m\times d$ 矩阵、$b\in\R^m$，则
 \[
  \varphi_{AX+b}(t)=e^{i\langle t,b\rangle}\varphi_X(A^Tt),
  \qquad t\in\R^m.
 \]
\end{enumerate}
\end{proposition}

\begin{proof}
对任意 $t,h\in\R^d$，
\[
 |\varphi_X(t+h)-\varphi_X(t)|
 \leq\E|e^{i\langle h,X\rangle}-1|.
\]
右端与 $t$ 无关，并由支配收敛定理在 $h\to0$ 时趋于零，所以连续性是一致的。
此外
\begin{align*}
 \sum_{j,k}c_j\overline{c_k}\varphi_X(t_j-t_k)
 &=\E\sum_{j,k}c_j\overline{c_k}
       e^{i\langle t_j-t_k,X\rangle}\\
 &=\E\left|\sum_jc_je^{i\langle t_j,X\rangle}\right|^2\geq0,
\end{align*}
这证明正定性。

若 $X,Y$ 独立，则期望因子化给
\[
 \varphi_{X+Y}(t)
 =\E\bigl(e^{i\langle t,X\rangle}e^{i\langle t,Y\rangle}\bigr)
 =\varphi_X(t)\varphi_Y(t).
\]
最后，内积恒等式
$\langle t,AX+b\rangle=\langle A^Tt,X\rangle+\langle t,b\rangle$
直接给出线性变换公式。
\end{proof}

第一项重要计算仍是高斯分布，但这次不能把实矩母函数公式中的 $t$ 形式地替换为 $it$；
那样做已经假定了尚未建立的复参数理论。下面从实变量积分重新计算。

\begin{example}[高斯特征函数]\label{ex:5-4-3-1}
令
\[
 F(t)=\frac1{\sqrt{2\pi}}\int_{\R}e^{itx}e^{-x^2/2}\dd x.
\]
对 $t$ 在任意有界区间变化时，导数被 $|x|e^{-x^2/2}$ 支配，故可在积分号下求导。
又 $xe^{-x^2/2}=-(e^{-x^2/2})'$。对 $R>0$ 在 $[-R,R]$ 上分部积分，
\begin{align*}
 \int_{-R}^{R}xe^{itx}e^{-x^2/2}\dd x
 &=-\int_{-R}^{R}e^{itx}\dd\bigl(e^{-x^2/2}\bigr)\\
 &=-\bigl[e^{itx}e^{-x^2/2}\bigr]_{-R}^{R}
   +it\int_{-R}^{R}e^{itx}e^{-x^2/2}\dd x.
\end{align*}
边界项的模不超过 $2e^{-R^2/2}$，因而趋于零；其余积分由
$e^{-x^2/2}$ 控制。令 $R\to\infty$ 得
\begin{align*}
 F'(t)
 &=\frac{i}{\sqrt{2\pi}}\int_{\R}xe^{itx}e^{-x^2/2}\dd x\\
 &=\frac{i}{\sqrt{2\pi}}\int_{\R}it e^{itx}e^{-x^2/2}\dd x
 =-tF(t).
\end{align*}
由 $F(0)=1$，实参数常微分方程的唯一解是
\[
 F(t)=e^{-t^2/2}.
\]

现在令 $Z=(Z_1,\dots,Z_m)$，其中坐标独立且均为 $N(0,1)$；对实
$d\times m$ 矩阵 $A$，暂把 $AZ$ 的分布记为 $N(0,\Sigma)$，其中
$\Sigma=AA^T$，允许 $\Sigma$ 奇异。独立乘法与线性变换公式给
\[
 \varphi_{AZ}(t)
 =\prod_{j=1}^m e^{-(A^Tt)_j^2/2}
 =\exp\left(-\frac12t^TAA^Tt\right)
 =\exp\left(-\frac12t^T\Sigma t\right).
\]
右端只含 $\Sigma$。定理~\ref{theorem:5-4-3-2} 将证明不同因子 $A$ 产生同一分布，
从而这个高斯分布记号不依赖所选分解。
\end{example}

如果两个概率的特征函数相同，我们知道它们对所有指数振荡测试相同；尚不能直接说“Fourier
反演”便结束，因为一般 Fourier 反演要到后章才建立。可用高斯核把两个测度先平滑，
而所需的高斯反演恒等式可以在本节用实变量常微分方程独立证明。

\begin{theorem}[特征函数唯一性定理]\label{theorem:5-4-3-2}
若 $\mu,\nu$ 是 $\R^d$ 上的概率测度，且
$\varphi_\mu(t)=\varphi_\nu(t)$ 对每个 $t\in\R^d$ 成立，则 $\mu=\nu$。
\end{theorem}

\begin{proof}
对 $\varepsilon>0$，记
\[
 g_\varepsilon(z)=(2\pi\varepsilon)^{-d/2}
 \exp\left(-\frac{|z|^2}{2\varepsilon}\right).
\]
我们先证明局部需要的高斯恒等式
\begin{equation}\label{eq:gaussian-local-inversion}
 g_\varepsilon(z)=\frac1{(2\pi)^d}
 \int_{\R^d}e^{-i\langle t,z\rangle}
 e^{-\varepsilon|t|^2/2}\dd t.
\end{equation}
把右端记为 $F_\varepsilon(z)$。高斯权重乘任意一次多项式仍可积，故可对每个
$z_j$ 在积分号下求导。对 $t_j$ 作实变量分部积分，并用
$\partial_{t_j}e^{-\varepsilon|t|^2/2}=-\varepsilon t_j
e^{-\varepsilon|t|^2/2}$，得到
\begin{align*}
 \partial_{z_j}F_\varepsilon(z)
 &=-\frac{i}{(2\pi)^d}\int t_j e^{-i\langle t,z\rangle}
 e^{-\varepsilon|t|^2/2}\dd t\\
 &=-\frac{z_j}{\varepsilon}F_\varepsilon(z).
\end{align*}
分部积分的边界项由高斯衰减趋于零。另一方面，多维实高斯积分的乘积分解给
\[
 F_\varepsilon(0)
 =\frac1{(2\pi)^d}\left(\frac{2\pi}{\varepsilon}\right)^{d/2}
 =(2\pi\varepsilon)^{-d/2}.
\]
沿每个坐标依次解上述实参数常微分方程，得到
$F_\varepsilon(z)=F_\varepsilon(0)e^{-|z|^2/(2\varepsilon)}$，即
\eqref{eq:gaussian-local-inversion}。

卷积 $\mu*g_\varepsilon$ 有密度
\begin{align*}
 p_{\mu,\varepsilon}(x)
 &=\int g_\varepsilon(x-y)\mu(\dd y)\\
 &=\frac1{(2\pi)^d}\int_{\R^d}
 e^{-i\langle t,x\rangle}e^{-\varepsilon|t|^2/2}
 \varphi_\mu(t)\dd t.
\end{align*}
Fubini 定理可用，因为高斯权重可积而 $|\varphi_\mu|\leq1$。
所以 $\varphi_\mu=\varphi_\nu$ 蕴含
$\mu*g_\varepsilon=\nu*g_\varepsilon$ 对每个 $\varepsilon>0$ 成立。

取一个分布为 $\mu$ 的随机向量 $X$，并在乘积空间上取与它独立的标准高斯向量 $Z$。
则 $X+\sqrt\varepsilon Z$ 的分布是 $\mu*g_\varepsilon$。对任意 $f\in C_b(\R^d)$，
$f(X+\sqrt\varepsilon Z)\to f(X)$ 逐点成立，且被 $\|f\|_\infty$ 支配，故
\[
 \int f\dd(\mu*g_\varepsilon)\longrightarrow\int f\dd\mu.
\]
同理，$\nu*g_\varepsilon\Rightarrow\nu$。由于两个平滑测度对每个 $\varepsilon$ 相同，
$\mu$ 与 $\nu$ 对所有 $C_b$ 函数积分相同，特别地对所有 $C_c$ 函数积分相同。
定理~\ref{theorem:3-1-1-4} 用 $C_c$ 测试恢复 Radon 开集质量，所以二者在所有开集、继而在
所有 Borel 集上相同。
\end{proof}

因此例~\ref{ex:5-4-3-1} 中若 $AA^T=BB^T$，则 $AZ$ 与 $BZ'$ 的特征函数相同，
唯一性定理保证它们同分布。奇异协方差也不例外；唯一性证明从未要求密度存在。

\begin{remark}[联合高斯分布的特征函数表示]\label{rem:5-4-3-joint-gaussian-recovery}
回到例~\ref{ex:5-2-3-1}中的联合高斯定义。设
$W=(W_1,\ldots,W_d)$ 满足该例的联合高斯定义，记
$m=\E W$ 且 $\Sigma=(\Cov(W_j,W_k))_{j,k}$。对任意 $t\in\R^d$，
$\langle t,W-m\rangle$ 或为常数零，或为均值零、方差
$t^T\Sigma t$ 的一维高斯变量。上面的一维计算因而给出
\[
 \varphi_W(t)=\exp\!\left(i\langle t,m\rangle-\frac12t^T\Sigma t\right).
\]
$\Sigma$ 半正定，故线性代数给出实矩阵 $A$ 使 $AA^T=\Sigma$。若
$Z$ 的坐标独立且均为 $N(0,1)$，则 $m+AZ$ 具有同一特征函数；由
定理~\ref{theorem:5-4-3-2}，
\[
 W\overset{\mathrm d}=m+AZ.
\]
因此“所有线性组合均为高斯分布”与上面的线性构造给出同一类分布。
若 $\Sigma$ 正定，在标准高斯乘积密度中作可逆线性换元，又得
\[
 p_W(x)=\frac{1}{(2\pi)^{d/2}(\det\Sigma)^{1/2}}
 \exp\!\left[-\frac12(x-m)^T\Sigma^{-1}(x-m)\right].
\]
$d=2$ 时这正是例~\ref{ex:5-2-3-1} 完成平方所用的联合密度。若
$\Sigma$ 奇异，分布集中于一个真仿射子空间，不应强行写成
$d$ 维 Lebesgue 密度；线性构造与特征函数仍然有效。
\end{remark}

特征函数比矩稳健得多。下一例用纯实变量卷积证明 Cauchy 分布的稳定性，同时避免使用留数定理。

\begin{example}[Cauchy 平均的稳定性]\label{ex:5-4-3-cauchy-stability}
对 $a>0$，令
\[
 f_a(x)=\frac{a}{\pi(a^2+x^2)}.
\]
我们证明 $f_a*f_b=f_{a+b}$。固定 $x$，置
\[
 \Delta=(x^2+(a+b)^2)(x^2+(a-b)^2).
\]
当 $\Delta\ne0$ 时，有部分分式分解
\begin{align*}
 \frac1{(y^2+b^2)((x-y)^2+a^2)}
 &=\frac{Ay+B}{y^2+b^2}
   +\frac{-A(y-x)+C}{(y-x)^2+a^2},
\end{align*}
其中
\[
 A=\frac{2x}{\Delta},\qquad
 B=\frac{x^2+a^2-b^2}{\Delta},\qquad
 C=\frac{x^2-a^2+b^2}{\Delta}.
\]
将上式两边乘以两个二次分母后，所需验证的多项式恒等式为
\[
 (Ay+B)\bigl((x-y)^2+a^2\bigr)
 +\bigl(-A(y-x)+C\bigr)(y^2+b^2)=1;
\]
左边的 $y^3,y^2,y$ 和常数项系数分别是
\begin{align*}
 &A-A=0,\\
 &-Ax+B+C,\\
 &A(x^2+a^2-b^2)-2Bx,\\
 &B(x^2+a^2)+(Ax+C)b^2.
\end{align*}
代入 $A,B,C$ 后，中间两项为
\[
 \frac{-2x^2+(x^2+a^2-b^2)+(x^2-a^2+b^2)}{\Delta}=0,
\]
以及
\[
 \frac{2x(x^2+a^2-b^2)-2x(x^2+a^2-b^2)}{\Delta}=0.
\]
常数项的分子则是
\begin{align*}
 &(x^2+a^2-b^2)(x^2+a^2)+(3x^2-a^2+b^2)b^2\\
 &=x^4+2(a^2+b^2)x^2+(a^2-b^2)^2
 =\Delta,
\end{align*}
故常数项等于一，多项式恒等式得证。
在线段 $[-R,R]$ 上积分时，两个线性分子项必须先合并。它们的原函数之和含有
\[
 \frac A2\log\frac{y^2+b^2}{(y-x)^2+a^2};
\]
先在两个端点作差、再令 $R\to\infty$，该对数贡献趋于零。不能把两个并不各自
Lebesgue 可积的线性分子项分开积分。常数项则给
\begin{align*}
 \int_{\R}\frac{\dd y}{(y^2+b^2)((x-y)^2+a^2)}
 &=\pi\left(\frac B b+\frac C a\right)\\
 &=\frac{\pi(a+b)}{ab\,[x^2+(a+b)^2]}.
\end{align*}
最后一步的代数化简是
\begin{align*}
 \frac Bb+\frac Ca
 &=\frac{a(x^2+a^2-b^2)+b(x^2-a^2+b^2)}{ab\Delta}\\
 &=\frac{(a+b)\bigl(x^2+(a-b)^2\bigr)}
 {ab\bigl(x^2+(a+b)^2\bigr)\bigl(x^2+(a-b)^2\bigr)}\\
 &=\frac{a+b}{ab\bigl(x^2+(a+b)^2\bigr)}.
\end{align*}
乘以 $ab/\pi^2$ 即得 $(f_a*f_b)(x)=f_{a+b}(x)$。唯一未被上述代数覆盖的情形是
$x=0$ 且 $a=b$；两边关于 $x$ 连续，故由 $x\to0$ 得到同一等式。
这里左边的连续性也可直接由 $f_a$ 的一致连续性估计：
\[
 |(f_a*f_b)(x+h)-(f_a*f_b)(x)|
 \leq\|\tau_hf_a-f_a\|_\infty\|f_b\|_1\longrightarrow0.
\]

所以独立的 Cauchy 变量相加时尺度参数相加。若 $X_1,\dots,X_n$ 独立且密度均为 $f_1$，
则 $S_n$ 的密度为 $f_n$。变量 $S_n/n$ 的密度为
\[
 y\longmapsto n f_n(ny)
 =\frac1{\pi(1+y^2)}=f_1(y).
\]
这证明 $S_n/n\overset{\mathrm d}=X_1$，也就完成了例~\ref{ex:5-4-2-2}所用的稳定性计算。
若另行证明 Fourier 恒等式 $\varphi_{f_a}(t)=e^{-a|t|}$，还可把卷积恒等式化成一次乘法；
本例不调用该恒等式，上面的推导完全建立在实变量积分上。
\end{example}

所有整数矩存在也未必能唯一识别分布；这不是一句抽象警告，而有一族可以逐项计算的密度。

\begin{remark}[所有整数矩仍可能不唯一决定分布]\label{rem:5-4-3-moment-indeterminate}
对 $x>0$ 令
\[
 p_0(x)=\frac1{x\sqrt{2\pi}}e^{-(\log x)^2/2},
 \qquad
 p_\theta(x)=p_0(x)\bigl[1+\theta\sin(2\pi\log x)\bigr],
 \quad |\theta|\leq1.
\]
每个 $p_\theta$ 非负。作换元 $y=\log x$。对整数 $n\geq0$ 和实数 $b$，配方并平移
$y=u+n$，再使用例~\ref{ex:5-4-3-1} 的高斯余弦积分，得到
\begin{align*}
 &\frac1{\sqrt{2\pi}}\int_{\R}
 e^{ny-y^2/2}\sin(by)\dd y\\
 &\qquad=e^{n^2/2}\frac1{\sqrt{2\pi}}
 \int_{\R}e^{-u^2/2}\bigl[\sin(bu)\cos(bn)+\cos(bu)\sin(bn)\bigr]\dd u\\
 &\qquad=e^{n^2/2}\left[0\cdot\cos(bn)+e^{-b^2/2}\sin(bn)\right]
 =e^{(n^2-b^2)/2}\sin(bn).
\end{align*}
其中第一项因被积函数为奇函数而为零。相同的配方还给出
\[
 \int_0^\infty x^np_0(x)\dd x
 =\frac1{\sqrt{2\pi}}\int_{\R}e^{ny-y^2/2}\dd y
 =e^{n^2/2}<\infty,
\]
而取 $n=0$ 时特别得到
\[
 \int_0^\infty p_0(x)\dd x
 =\frac1{\sqrt{2\pi}}\int_{\R}e^{-y^2/2}\dd y=1.
\]
取 $b=2\pi$，便有
\[
 \int_0^\infty x^np_0(x)\sin(2\pi\log x)\dd x
 =e^{(n^2-4\pi^2)/2}\sin(2\pi n)=0.
\]
结合上式，$n=0$ 的扰动积分为零说明每个 $p_\theta$ 都归一化；一般 $n$ 说明它们的
$n$ 阶矩都存在且等于 $e^{n^2/2}$。
然而当 $\theta\ne\theta'$ 时，两个密度在
$\sin(2\pi\log x)\ne0$ 的正测度集合上不同。因此“全部矩存在”本身不足以决定分布。
这不排除在额外矩增长条件下成立的矩唯一性定理。
\end{remark}


点态极限还剩一个危险：它可能不再是特征函数。最简单的失败发生在质量整体逃向无穷时。

\begin{example}[Bernoulli 与 Rademacher 的零点展开]\label{ex:5-4-3-2}
若 $R$ 为 Rademacher 变量，则 $\varphi_R(t)=\cos t$，所以
\[
 \varphi_R(t)=1-\frac{t^2}{2}+o(t^2)\qquad(t\to0).
\]
对独立副本 $R_i$，
\[
 \varphi_{n^{-1/2}\sum_{i=1}^nR_i}(t)
 =\cos(t/\sqrt n)^n\longrightarrow e^{-t^2/2}.
\]
更一般地，若 $B\sim\operatorname{Bernoulli}(p)$、$0<p<1$，则
$X=(B-p)/\sqrt{p(1-p)}$。记 $q=1-p$，则 $X$ 以概率 $p$ 取
$\sqrt{q/p}$，以概率 $q$ 取 $-\sqrt{p/q}$，所以
\begin{align*}
 \varphi_X(t)
 &=p e^{it\sqrt{q/p}}+q e^{-it\sqrt{p/q}}\\
 &=p\left(1+it\sqrt{\frac qp}-\frac{t^2q}{2p}+o(t^2)\right)
  +q\left(1-it\sqrt{\frac pq}-\frac{t^2p}{2q}+o(t^2)\right)\\
 &=1+it(\sqrt{pq}-\sqrt{pq})-\frac{t^2}{2}(q+p)+o(t^2)\\
 &=1-\frac{t^2}{2}+o(t^2).
\end{align*}
因而它与 Rademacher 变量具有同一个二阶主项；偏斜度只进入更高阶余项。
\end{example}


\begin{example}[零点处不连续的极限]\label{ex:5-4-3-3}
令 $\mu_n$ 为 $[-n,n]$ 上的均匀概率。直接积分给
\[
 \varphi_n(t)=
 \begin{cases}
  \dfrac{\sin(nt)}{nt},&t\ne0,\\[5pt]
  1,&t=0.
 \end{cases}
\]
对每个 $t\ne0$，$\varphi_n(t)\to0$，而在 $t=0$ 处始终等于一，所以点态极限在零点
不连续。另一方面，对固定 $R>0$ 与 $n\geq R$，
$\mu_n([-R,R])=R/n\to0$；任意紧集都包含在某个这样的区间中，因而没有一个紧集能统一抓住
这些概率的大部分质量。零点连续性将在下面的证明中恰好变成紧性。
\end{example}

\begin{theorem}[\Levy{} 连续性定理]\label{theorem:5-4-3-3}
设 $\mu_n,\mu$ 是 $\R^d$ 上的概率测度。
\begin{enumerate}
 \item 若 $\mu_n\Rightarrow\mu$，则
 $\varphi_{\mu_n}(t)\to\varphi_\mu(t)$ 对每个 $t\in\R^d$ 成立；
 \item 反之，若 $\varphi_{\mu_n}(t)\to\varphi(t)$ 对每个 $t$ 成立，且
 $\varphi$ 在 $0$ 处连续，则存在概率测度 $\mu$，使
 $\varphi=\varphi_\mu$ 且 $\mu_n\Rightarrow\mu$。
\end{enumerate}
\end{theorem}

\begin{proof}
若 $\mu_n\Rightarrow\mu$，对固定 $t$，函数
$x\mapsto e^{i\langle t,x\rangle}$ 有界连续，故弱收敛的定义直接给第一部分。

现在假设第二部分的条件。先证 $(\mu_n)$ 是紧族。先处理 $d=1$，并在规范空间上令
$X_n(x)=x$、$X_n\sim\mu_n$。对 $T>0$，Fubini 定理给
\begin{align*}
 \frac1{2T}\int_{-T}^T\Re\varphi_{\mu_n}(t)\dd t
 &=\E\left[\frac1{2T}\int_{-T}^T\cos(tX_n)\dd t\right]\\
 &=\E\left[\frac{\sin(TX_n)}{TX_n}\right],
\end{align*}
其中 $X_n=0$ 时把比值定义为一。若 $|X_n|\geq2/T$，则
$|\sin(TX_n)/(TX_n)|\leq1/2$，所以
\[
 \1_{\{|X_n|\geq2/T\}}
 \leq2\left(1-\frac{\sin(TX_n)}{TX_n}\right).
\]
取期望得到
\begin{equation}\label{eq:levy-tightness-estimate}
 \Prob(|X_n|\geq2/T)
 \leq2\left(1-\frac1{2T}\int_{-T}^T
                 \Re\varphi_{\mu_n}(t)\dd t\right).
\end{equation}
固定 $T$ 后，由 $|\varphi_{\mu_n}|\leq1$、点态收敛和支配收敛定理，右端收敛到
\[
 2\left(1-\frac1{2T}\int_{-T}^T\Re\varphi(t)\dd t\right).
\]
$\varphi(0)=\lim_n\varphi_{\mu_n}(0)=1$，且它在零点连续，所以这个量在
$T\downarrow0$ 时趋于零。给定
$\varepsilon>0$，先取 $T$ 足够小，再由上述收敛取 $N$，便能用区间
$[-2/T,2/T]$ 抓住所有 $n\geq N$ 的至少 $1-\varepsilon$ 质量。由
\cref{proposition:5-3-3-polish-radon}，剩余有限多个概率各自可由一个
有界闭区间抓住至少 $1-\varepsilon$ 的质量；取这些区间与前一区间的有限并，得到全列的紧性。

在 $\R^d$ 中，对每个坐标轴方向 $e_j$ 把同一估计应用于
$t\mapsto\varphi_{\mu_n}(te_j)$。给定 $\varepsilon>0$，利用 $\varphi$ 在零点的连续性，
可取同一个充分小的 $T>0$，使每个坐标的极限上界都小于 $\varepsilon/(2d)$；再取共同的
$N$，使对所有 $n\geq N$、所有 $1\leq j\leq d$ 都有
\[
 \mu_n\bigl(\{x:|x_j|\geq2/T\}\bigr)<\varepsilon/d.
\]
并集界说明立方体 $[-2/T,2/T]^d$ 抓住这些尾项至少 $1-\varepsilon$ 的质量。剩余有限多个
概率各自具有紧逼近，把相应有限多个紧集与该立方体作有限并，便得到对全列有效的紧集。因此一般维数的
$(\mu_n)$ 也是紧族。

由 Prokhorov 定理~\ref{theorem:5-3-3-2}，存在子列
$\mu_{n_k}\Rightarrow\mu$。第一部分说明
\[
 \varphi_\mu(t)=\lim_k\varphi_{\mu_{n_k}}(t)=\varphi(t),
\]
所以 $\varphi$ 确为一个概率的特征函数。任取原序列的任一子列；其紧性不变，故仍有进一步
弱收敛子列，而其极限的特征函数仍为 $\varphi$。唯一性定理
\ref{theorem:5-4-3-2} 说明所有这类子列极限都等于 $\mu$。

最后说明整列收敛。若 $\mu_n$ 不弱收敛到 $\mu$，由命题
\ref{proposition:5-3-3-weak-metric} 的有界 Lipschitz 度量化，存在 $\eta>0$ 和一个子列
满足 $d_{\mathrm{BL}}(\mu_{n_k},\mu)\geq\eta$。该子列又有进一步弱收敛子列，且其极限必须是
$\mu$，这与 $d_{\mathrm{BL}}$ 沿该进一步子列趋于零矛盾。因此 $\mu_n\Rightarrow\mu$。
\end{proof}

证明把两项工作严格分开：Prokhorov 从紧性产生候选极限，特征函数唯一性识别所有候选。
只有唯一性而没有存在性，不能排除质量逃逸；例~\ref{ex:5-4-3-3} 正是这个缺口的可计算模型。

现在回到独立和。二阶矩在零点给出特征函数的二阶展开，而 $\sqrt n$ 缩放恰好让
$n$ 个二次项累计成有限指数。

\begin{definition}[标准化和]\label{def:5-4-3-3}
若 $(X_i)$ i.i.d.，$\E X_1=\mu$ 且
$0<\Var(X_1)=\sigma^2<\infty$，定义
\[
 S_n=\sum_{i=1}^nX_i,
 \qquad
 Z_n=\frac{S_n-n\mu}{\sigma\sqrt n}.
\]
\end{definition}

\begin{theorem}[Lindeberg--\Levy{} 中心极限定理]\label{theorem:5-4-3-4}
若 $(X_i)$ i.i.d.，$\E X_1=\mu$ 且
$0<\Var(X_1)=\sigma^2<\infty$，则
\[
 \frac{S_n-n\mu}{\sigma\sqrt n}\Rightarrow N(0,1).
\]
\end{theorem}

\begin{proof}
用 $Y_i=(X_i-\mu)/\sigma$ 替换 $X_i$，只需处理
$\E Y_1=0$、$\E Y_1^2=1$ 的情形。记 $\varphi(u)=\E e^{iuY_1}$。
对实数 $v$，积分余项公式给
\[
 e^{iv}-1-iv=-\int_0^v(v-s)e^{is}\dd s,
 \qquad
 |e^{iv}-1-iv|\leq\frac{v^2}{2}.
\]
把
\[
 r(v)=
 \begin{cases}
 \dfrac{e^{iv}-1-iv}{v^2},&v\ne0,\\[5pt]
 -\dfrac12,&v=0
 \end{cases}
\]
看作连续函数。事实上，在余项积分中作 $s=uv$ 换元，对每个 $v\in\R$ 都有
\[
 r(v)=-\int_0^1(1-u)e^{iuv}\dd u.
\]
该表示式直接给出 $r(v)\to-1/2$ 当 $v\to0$，并且
$|r(v)|\le\int_0^1(1-u)\dd u=1/2$。于是
\[
 \frac{e^{iuY_1}-1-iuY_1}{u^2}=Y_1^2r(uY_1).
\]
右端被 $Y_1^2/2$ 支配，并在 $u\to0$ 时趋于 $-Y_1^2/2$。
由支配收敛定理及 $\E Y_1=0$、$\E Y_1^2=1$，
\[
 \varphi(u)=1-\frac{u^2}{2}+o(u^2).
\]

独立性给出，对每个固定 $t\in\R$，
\[
 \varphi_{n^{-1/2}\sum_{i=1}^nY_i}(t)
 =\varphi(t/\sqrt n)^n.
\]
令 $w_n=n(\varphi(t/\sqrt n)-1)$，则 $w_n\to-t^2/2$，所以右端为
$(1+w_n/n)^n$。这里不需要选择复对数分支。事实上，对任意复数列 $w_n\to w$，
二项式展开写成
\[
 \left(1+\frac{w_n}{n}\right)^n
 =\sum_{k=0}^{\infty}
 \1_{\{k\leq n\}}\frac{n(n-1)\cdots(n-k+1)}{n^k}
 \frac{w_n^k}{k!}.
\]
对每个固定 $k$，通项趋于 $w^k/k!$；当 $n$ 充分大时，绝对值被
$M^k/k!$ 支配，其中 $M>\sup_n|w_n|$。在计数测度上使用支配收敛定理，得到极限
$\sum_{k\geq0}w^k/k!=e^w$。因此
\[
 \varphi(t/\sqrt n)^n\longrightarrow e^{-t^2/2}.
\]
例~\ref{ex:5-4-3-1} 已证明右端是 $N(0,1)$ 的特征函数，并且它在零点连续。
\Levy{} 连续性定理~\ref{theorem:5-4-3-3} 于是给出所需弱收敛。
\end{proof}

\begin{center}
\begin{tikzpicture}[node distance=8mm and 8mm]
 \node[book node] (local) {单项零点展开\\$\varphi(u)=1-u^2/2+o(u^2)$};
 \node[book strong node,right=of local] (scale) {$\sqrt n$ 缩放\\$u=t/\sqrt n$};
 \node[book node,right=of scale] (power) {独立乘积\\$\varphi(t/\sqrt n)^n$};
 \node[book warm node,below=of power] (gauss) {高斯极限\\$e^{-t^2/2}$};
 \node[book strong node,left=of gauss] (levy) {\Levy{} 连续性\\$\Rightarrow$ 分布收敛};
 \draw[book accent arrow] (local)--(scale);
 \draw[book accent arrow] (scale)--(power);
 \draw[book accent arrow] (power)--(gauss);
 \draw[book accent arrow] (gauss)--(levy);
\end{tikzpicture}
\end{center}

这张图也标出定理的边界。有限二阶矩用于二阶展开；独立性用于乘法；零点连续性与
\Levy{} 定理用于把函数极限升级为概率极限。定理本身不给收敛速度；要得到定量误差率，还需要
三阶绝对矩等附加假设。多维中心极限定理可沿同一路线对每个 $t\in\R^d$ 展开
$\langle t,X\rangle$；其陈述还须明确向量矩条件、协方差矩阵以及退化协方差的情形。

本小节还澄清了矩母函数与特征函数的分工。前者在存在时能给强尾界，却可能对重尾分布失效；
后者始终存在并唯一决定分布，却不自动提供指数尾部。Fourier 反演、复解析延拓和定量局部极限定理
都需要额外工具；这里的唯一性、\Levy{} 定理与中心极限定理证明只使用实变量积分、紧性和
特征函数的基本性质。

\begin{exercise}[特征函数与极限]
\begin{enumerate}
 \item 设 $N$ 取值于 $\{0,1,2,\ldots\}$，且
 $\Prob(N=k)=e^{-\lambda}\lambda^k/k!$（$k\geq0$，$\lambda>0$）。计算 $N$ 的特征函数，
 并由独立乘法验证两个此类独立变量之和仍属于同一分布族，其参数为原参数之和；
 \item 证明若 $X_n\Rightarrow X$ 且 $\sup_n\E|X_n|^2<\infty$，则 $(\Law(X_n))$ 是紧族；
       再说明这个矩界本身不能推出二阶矩收敛；
 \item 对 $B\sim\operatorname{Bernoulli}(p)$ 直接展开标准化变量的特征函数到二阶，并与
       定理~\ref{theorem:5-4-3-4} 对照；
 \item 检查例~\ref{ex:5-4-3-3} 的点态极限为何不能是任何概率的特征函数。
\end{enumerate}
\end{exercise}

至此，独立和的三种尺度已经分开：不缩放时集中不等式给有限样本尾部，除以 $n$ 后大数律给
确定极限，除以 $\sqrt n$ 后中心极限定理保留高斯波动。下一章把单个随机变量换成时间索引族；
有限维分布仍是有限观测，而一致性条件将决定这些有限观测能否拼成路径空间上的分布。
